\documentclass[12pt, a4paper]{article}

\usepackage[utf8]{inputenc}
\usepackage[T1]{fontenc}
\usepackage[french]{babel}
\usepackage{geometry}
\usepackage{booktabs}
\usepackage{array}
\usepackage{xcolor}
\usepackage{listings}
\usepackage{hyperref}
\usepackage{graphicx}
\usepackage{amsmath}
\usepackage{titlesec}
\usepackage{enumitem}
\usepackage{fancyhdr}

\geometry{margin=2.5cm}

\pagestyle{fancy}
\fancyhf{}
\rhead{Classification d'émotions -- Tweets COVID-19}
\lhead{\leftmark}
\rfoot{Page \thepage}

\definecolor{codeblue}{rgb}{0.13, 0.29, 0.53}
\definecolor{codegray}{rgb}{0.95, 0.95, 0.95}

\lstset{
    backgroundcolor=\color{codegray},
    basicstyle=\ttfamily\small,
    keywordstyle=\color{codeblue}\bfseries,
    breaklines=true,
    frame=single,
    language=Python,
    showstringspaces=false,
}

\hypersetup{
    colorlinks=true,
    linkcolor=codeblue,
    urlcolor=codeblue,
}

% ============================================================
\title{
    \textbf{Classification d'émotions dans des tweets COVID-19} \\[0.5em]
    \large Prétraitement, lemmatisation et régression logistique
}
\author{Kaufmann Raphaël}
\date{Juin 2026}
% ============================================================

\begin{document}

\maketitle
\tableofcontents
\newpage

% ------------------------------------------------------------
\section{Introduction}
% ------------------------------------------------------------

Ce document récapitule les étapes réalisées pour construire un pipeline de classification
d'émotions appliqué à des tweets relatifs au COVID-19.
L'objectif est d'associer à chaque tweet l'une des sept émotions suivantes :
\textbf{neutral, disgust, happiness, anger, surprise, sadness, fear}.

Le pipeline se décompose en trois grandes phases :
\begin{enumerate}
    \item \textbf{Prétraitement et filtrage} -- nettoyage du texte et extraction de features émotionnelles ;
    \item \textbf{Lemmatisation} -- réduction des mots à leur forme canonique, avec exceptions pour les mots d'emphase ;
    \item \textbf{Classification par régression logistique} -- entraînement et évaluation du modèle.
\end{enumerate}

% ------------------------------------------------------------
\section{Prétraitement des tweets}
% ------------------------------------------------------------

\subsection{Données source}

Le dataset initial est \texttt{covid19\_tweets.csv}, contenant environ 179\,000 tweets.
La colonne d'intérêt est \texttt{text}.
Un premier nettoyage effectué dans \texttt{indexation.py} supprime les URLs à l'aide
d'une expression régulière :

\begin{lstlisting}
motif_lien = r"https?://\S+"
tweet_nettoye = re.sub(motif_lien, '', tweet)
\end{lstlisting}

\subsection{Filtrage des mots courants (liste-vide)}

Le filtrage est réalisé dans \texttt{lemmatisation.py} via le modèle spaCy
\texttt{en\_core\_web\_sm}. Trois règles guident la conservation d'un token :

\begin{enumerate}
    \item \textbf{Mot utile} : le token n'appartient pas à la liste-vide spaCy (326 stop-words anglais).
    \item \textbf{Négation} : certains stop-words sont explicitement conservés car porteurs de sens
          émotionnel (colère, peur, tristesse).
    \item \textbf{Accentuation par répétition consécutive} : un stop-word répété $\geq 2$ fois
          d'affilée (ex. \textit{very very}) est conservé pour signaler une intensité émotionnelle.
\end{enumerate}

Les négations préservées sont :

\begin{center}
\texttt{not, no, nor, never, nothing, nobody, none, neither, nowhere, n't, cannot, without, most}
\end{center}

\subsection{Features émotionnelles extraites}

En parallèle du filtrage textuel, trois indicateurs numériques sont calculés
\textbf{sur le tweet original} (avant mise en minuscules) :

\begin{table}[h!]
\centering
\begin{tabular}{lll}
\toprule
\textbf{Feature} & \textbf{Description} & \textbf{Signal émotionnel} \\
\midrule
\texttt{pct\_majuscules} & \% de lettres majuscules & Énervement, intensité \\
\texttt{exclamations}    & Nombre de \texttt{!}     & Colère, surprise, joie forte \\
\texttt{interrogations}  & Nombre de \texttt{?}     & Confusion, incrédulité \\
\bottomrule
\end{tabular}
\caption{Features numériques extraites pour la détection d'émotions}
\end{table}

\subsection{Conservation des émojis}

Les émojis sont détectés via la librairie \texttt{emoji} et systématiquement conservés,
car ils constituent des marqueurs d'émotion particulièrement explicites
(\texttt{😡} → colère, \texttt{😢} → tristesse).

% ------------------------------------------------------------
\section{Lemmatisation}
% ------------------------------------------------------------

\subsection{Principe}

La lemmatisation réduit chaque mot à sa forme de base :
\textit{running} $\rightarrow$ \textit{run}, \textit{deaths} $\rightarrow$ \textit{death}.
Elle est appliquée après le filtrage, via le pipeline spaCy avec NER et parser désactivés
pour des raisons de performance (\texttt{nlp.pipe(..., disable=["ner", "parser"])}).

\subsection{Exception : mots d'emphase émotionnelle}

Les mots comportant une lettre répétée \textbf{3 fois ou plus consécutivement} ne sont
\textbf{pas lemmatisés}, car cette répétition de lettres est elle-même un signal d'émotion :

\begin{center}
\textit{WHATTTTT}, \textit{heyyyy}, \textit{nooooo}
\end{center}

La détection repose sur l'expression régulière \texttt{(.)\textbackslash1\{2,\}} :
si un caractère apparaît au moins 3 fois de suite, le mot est conservé tel quel.

\subsection{Format de sortie}

Chaque tweet traité est exporté en chaîne de caractères (mots séparés par des espaces),
puis sauvegardé dans \texttt{data/tweets\_traites.csv} avec les colonnes suivantes :

\begin{center}
\texttt{tweet\_lemmatise ;\ pct\_majuscules ;\ exclamations ;\ interrogations}
\end{center}

% ------------------------------------------------------------
\section{Classification par régression logistique}
% ------------------------------------------------------------

\subsection{Construction des features}

Le modèle combine deux familles de features via un \texttt{ColumnTransformer} scikit-learn :

\begin{itemize}
    \item \textbf{TF-IDF} sur le texte : unigrammes + bigrammes, 5\,000 features max,
          fréquence minimale de 2 occurrences.
    \item \textbf{Features numériques} : \texttt{pct\_majuscules}, \texttt{exclamations},
          \texttt{interrogations}, normalisées par \texttt{StandardScaler}.
\end{itemize}

\subsection{Modèle}

\begin{lstlisting}
LogisticRegression(
    max_iter=1000,
    class_weight="balanced",  # compense le desequilibre des classes
    C=1.0,
    random_state=42,
)
\end{lstlisting}

Le paramètre \texttt{class\_weight="balanced"} est crucial : il pondère automatiquement
les classes sous-représentées pour éviter que \textit{neutral} (majoritaire) écrase les autres.

\subsection{Première expérience : 193 tweets étiquetés manuellement}

\begin{table}[h!]
\centering
\begin{tabular}{lrrrr}
\toprule
\textbf{Émotion} & \textbf{Exemples} & \textbf{Précision} & \textbf{Rappel} & \textbf{F1} \\
\midrule
neutral   & 111 & 0.40 & 0.30 & 0.34 \\
fear      &  52 & 0.56 & 0.36 & 0.43 \\
happiness &  24 & 0.00 & 0.00 & 0.00 \\
sadness   &   3 & 0.00 & 0.00 & 0.00 \\
anger     &   2 & 0.00 & 0.00 & 0.00 \\
disgust   &   1 & 0.00 & 0.00 & 0.00 \\
\midrule
\multicolumn{2}{l}{\textbf{Accuracy globale}} & \multicolumn{3}{c}{\textbf{28\,\%}} \\
\bottomrule
\end{tabular}
\caption{Résultats avec 193 tweets étiquetés -- dataset trop petit et déséquilibré}
\end{table}

Les classes avec moins de 5 exemples ne peuvent pas être apprises correctement.

\subsection{Étiquetage automatique : passage à 3\,000 tweets}

Pour pallier le manque de données, un modèle pré-entraîné de HuggingFace a été utilisé
pour étiqueter automatiquement 3\,000 tweets :

\begin{center}
\texttt{j-hartmann/emotion-english-distilroberta-base}
\end{center}

Ce modèle, basé sur DistilRoBERTa, détecte 7 émotions alignées avec nos cibles.
L'émotion \textit{joy} est remappée vers \textit{happiness}.

\begin{table}[h!]
\centering
\begin{tabular}{lr}
\toprule
\textbf{Émotion} & \textbf{Tweets étiquetés} \\
\midrule
neutral   & 1613 \\
fear      &  378 \\
sadness   &  357 \\
happiness &  298 \\
anger     &  220 \\
surprise  &  119 \\
disgust   &   15 \\
\midrule
\textbf{Total} & \textbf{3\,000} \\
\bottomrule
\end{tabular}
\caption{Répartition des émotions après étiquetage automatique}
\end{table}

\subsection{Deuxième expérience : 3\,000 tweets auto-étiquetés}

\begin{table}[h!]
\centering
\begin{tabular}{lrrrr}
\toprule
\textbf{Émotion} & \textbf{Support} & \textbf{Précision} & \textbf{Rappel} & \textbf{F1} \\
\midrule
neutral   & 323 & 0.76 & 0.52 & 0.62 \\
anger     &  44 & 0.42 & 0.57 & 0.49 \\
fear      &  76 & 0.41 & 0.50 & 0.45 \\
sadness   &  71 & 0.39 & 0.51 & 0.44 \\
happiness &  59 & 0.30 & 0.46 & 0.36 \\
surprise  &  24 & 0.14 & 0.25 & 0.18 \\
disgust   &   3 & 0.00 & 0.00 & 0.00 \\
\midrule
\multicolumn{2}{l}{\textbf{Accuracy globale}} & \multicolumn{3}{c}{\textbf{50\,\%}} \\
\bottomrule
\end{tabular}
\caption{Résultats avec 3\,000 tweets -- nette amélioration}
\end{table}

L'accuracy passe de \textbf{28\,\%} à \textbf{50\,\%}, soit un gain de 22 points.
Les mots les plus prédictifs par émotion confirment la cohérence du modèle :

\begin{itemize}
    \item \textbf{anger} : \textit{protest, blame, hate, refuse, stupidity}
    \item \textbf{fear} : \textit{lockdown, risk, crisis, quarantine, threat}
    \item \textbf{happiness} : \textit{good, happy, thank, great, love}
    \item \textbf{sadness} : \textit{death, miss, suffer, depression, bad}
    \item \textbf{neutral} : \textit{update, mask, case, test, community}
\end{itemize}

% ------------------------------------------------------------
\section{Conclusion et pistes d'amélioration}
% ------------------------------------------------------------

Le pipeline construit permet de classifier automatiquement les émotions de tweets COVID-19
avec une accuracy de 50\,\% à partir de 3\,000 exemples.

Les principales limites identifiées sont :
\begin{itemize}
    \item \textbf{Déséquilibre persistant} : \textit{disgust} reste quasi absent (15 exemples).
    \item \textbf{Qualité des labels} : l'étiquetage automatique introduit du bruit.
    \item \textbf{Modèle simple} : la régression logistique ne capture pas les relations
          syntaxiques complexes.
\end{itemize}

Pistes pour améliorer les performances :
\begin{enumerate}
    \item Augmenter le dataset à 10\,000+ tweets étiquetés.
    \item Fine-tuner directement un modèle BERT/RoBERTa sur les tweets COVID.
    \item Valider manuellement un sous-ensemble des labels automatiques.
\end{enumerate}

\end{document}